%
% Complete documentation on the extended LaTeX markup used for Insight
% documentation is available in ``Documenting Insight'', which is part
% of the standard documentation for Insight.  It may be found online
% at:
%
%     http://www.itk.org/

\documentclass{InsightArticle}


%%%%%%%%%%%%%%%%%%%%%%%%%%%%%%%%%%%%%%%%%%%%%%%%%%%%%%%%%%%%%%%%%%
%
%  hyperref should be the last package to be loaded.
%
%%%%%%%%%%%%%%%%%%%%%%%%%%%%%%%%%%%%%%%%%%%%%%%%%%%%%%%%%%%%%%%%%%
\usepackage[dvips,
bookmarks,
bookmarksopen,
backref,
colorlinks,linkcolor={blue},citecolor={blue},urlcolor={blue},
]{hyperref}
% to be able to use options in graphics
\usepackage{graphicx}
% for pseudo code
\usepackage{listings}
% subfigures
\usepackage{subfigure}


%  This is a template for Papers to the Insight Journal. 
%  It is comparable to a technical report format.

% The title should be descriptive enough for people to be able to find
% the relevant document. 
\title{Efficient implementation of kernel filtering}

% Increment the release number whenever significant changes are made.
% The author and/or editor can define 'significant' however they like.
\release{0.00}

% At minimum, give your name and an email address.  You can include a
% snail-mail address if you like.
\author{Richard Beare{$^1$} {\small{and}} Ga\"etan Lehmann{$^2$}}
\authoraddress{{$^1$}Department of Medicine, Monash University, Australia.\\
{$^2$}INRA, UMR 1198; ENVA; CNRS, FRE 2857, Biologie du D\'eveloppement et
Reproduction, Jouy en Josas, F-78350, France}

\begin{document}
\maketitle

\ifhtml
\chapter*{Front Matter\label{front}}
\fi


\begin{abstract}
\noindent
% The abstract should be a paragraph or two long, and describe the
% scope of the document.
Kernel based filtering is one of the fundamental tools of image
analysis and processing. A number of approaches have been developed
over the years that allow efficient implementation of such filters
even when the kernel size is large. This article reviews some of these
methods and introduces their ITK implementations.
\end{abstract}

\tableofcontents

\section{Introduction}
A kernel based filtering process replaces the pixel at the kernel
origin with the result of applying a function to all pixels defined by
the kernel. Many useful filters, including edge detection and gradient
filters, smoothing filters and rank and morphology filters fall into
this category. Direct implementations of such filters typically
involve visiting all pixels defined by the kernel in order to evaluate
the filter function. Such an approach is usually easy to implement
-- in ITK it is made simple by the neighborhood iterators -- but leads
to an algorithm complexity proportional to the number of pixels in the
kernel (or $O(n^d)$, where $n$ is the kernel size and $d$ is the
dimensionality). Such complexity tends to restrict application of such
filters to small kernels.

A number of classical methods exist for reducing this complexity to
more manageable, in some cases kernel size independent, levels. ITK
already exploits such techniques for Gaussian convolution
operations. This paper describes the classical techniques for
optimized mathematical morphology filters, rank filters and certain
convolution filters. These filters make the use of large kernels,
which can be very useful in many applications, practical on
conventional computing hardware.

\section{Separability and recursive implementations}
The two approaches most typically used to reduce complexity of kernel
based filters are separability and recursive computation, both of
which are used in the ITK implementation of Gaussian convolution
filters. A separable filter implements a multidimensional kernel by
cascading several one dimensional kernels, therefore reducing
complexity from $O(n^d)$ to $O(nd)$. The second approach exploits
redundancy that might be present in the computations of kernel
functions at neighboring locations, leading, in some cases, to a complexity
independent of $n$.

\section{Mathematical morphology operations}
Two optimized forms of the mathematical morphology operations of
erosion, dilation, opening and closing are presented here. The kernel
is usually referred to as the structuring element in mathematical
morphology. The methods described here are applied to ``flat''
structuring elements, that is structuring elements without
weights. The first method can be applied to arbitrary structuring
elements while the second can be applied to line structuring elements.

\subsection{Arbitrary structuring elements}
\label{sect:MMmovingHist}
The method described in \cite{Vandroogenbroeck96.3} relies on the
simple concept of an updatable histogram, and is often described as a
``moving histogram'' approach. A histogram is computed for a kernel
located at the first voxel. The histogram at the neighboring voxel can
then be computed by including newly included voxels and removing newly
excluded voxels. The list of included and excluded voxels
corresponding to movement in any direction can be computed when the
structuring element is created, and the direction with the smallest
number of changes should be selected as the direction for sweeping the
kernel across the image. The erosion or dilation at each location is
computed by selecting the minimum or maximum from the histogram. This
approach is very efficient when 8 or 16 bit pixels are used because
the histogram can be represented as an array, with place holders
tracking the current maximum or minimum increasing performance. More
sophisticated histogram representations are necessary for larger pixel
types. Our implementation uses c++ maps. This is an extension to the
original paper.

This methodology reduces the complexity from $O(n^d)$ to $O(n^{d-1})$,
while keeping the structuring element identical to the direct
implementation.

\subsection{Decomposition of structuring elements}
\label{sect:MMdecomp}
Morphological erosions and dilations are separable -- successive
dilation by orthogonal lines is equivalent to dilation by a rectangle
with sides equal to the line lengths. This means that any
hyper-rectangular structuring element can be constructed using lines
parallel to the image axes.

Approximations of more complex shapes, particularly circles and
ellipses can constructed using a number of lines at evenly spaced
angles \cite{Adams93}. It is difficult to create a structuring element
with a precisely defined radius using this method because line
structuring elements from which the circle structuring element is
composed must have odd length and there are practical limits due to
the realities of underlying digital grid representation of images. In
addition it is possible that the structuring elements may not be truly
translation invariant due to the representation of line
(e.g. Bresenham) used in the decomposition. However, precisely defined
radii are rarely critical when a large structuring element is called
for. An example of a structuring element created using line
structuring elements is shown in Figure
\ref{fig:circledecomposition}. Composition of regular shapes, such as
hexagons and octagons is more accurate.

\begin{figure}[htbp]
\centering
\includegraphics{kernel}
\caption{Approximate circular structuring element, radius 25, constructed using 8 lines.\label{fig:circledecomposition}}
\end{figure}

It is also possible, in theory, to construct 3D structuring elements
in similar ways. The construction of a hyper-rectangle is trivial,
however construction of spheres is more problematic. The code
discussed later provides preliminary implementations based on some
platonic solids and various spherical approximations, but further
testing and development is needed.

\subsection{Line structuring elements}
The decompositions discussed above are important because an efficient,
recursive, implementation of erosion and dilations along lines
exists. This method was introduced in \cite{Gil1993,vanHerk1992a} and
can compute an erosion or dilation in 3 operations per pixel,
independent of structuring element length. \cite{Gil2000} recently
reduced the cost to 1.5 pixels per pixel, but the procedure is more
complicated and there are reports of no speedup in reality.

The original algorithms utilize forward and backward running maxima
(for dilation) for the length of the structuring element from a pixel
of interest. The dilation can then be computed for a region the size
of the structuring element around the point of interest by comparing
values on the running extrema that are separated by the structuring
element length. This is illustrated in Figure \ref{fig:vHGWmethod}

\begin{figure}[htbp]
\centering
\includegraphics[scale=0.5]{vHGWexpl}
\caption{The van Herk, Gil, Werman method.\label{fig:vHGWmethod}}
\end{figure}

A method published recently that offers improved performance
\cite{Vandroogenbroeck2005Morphological} and the ability to perform a
line opening directly (rather than using an erosion/dilation
cascade). This method, known as {\em anchor morphology} has been
implemented in the filters discussed later in the article, but
performance issues are not yet clear. The method employs histograms
and therefore needs more complex data structures when applied to
higher precision data, potentially reducing any speed advantage.

\section{Rank filters}
\label{sect:rank}
Efficient implementations of median and rank filters can be carried out
using exactly the same approach as discussed for morphological
operations with arbitrary structuring elements. The method was
originally proposed in \cite{Huang79}. The implementation discussed
later supports arbitrary kernel shapes and pixel types as well as any
choice of rank.

Rank filters are not separable. However the performance benefits
offered by separability make it worth pretending they are. If median
filtering is being used to provide robust noise filtering or
background estimation then a separable approximation is worth
testing. This concept was originally proposed in \cite{Narendra81}. In
fact it is also possible to apply the decomposition discussed in
Section \ref{sect:MMdecomp} to median filters to achieve a closer
approximation to circular kernels.

\section{Mean and variance filters}
\label{sect:meanVar}
A simplified version of the moving histogram approach can also be used
to efficiently compute the mean or variance of an arbitrarily shaped
kernel. The histogram used for morphology and rank filters is a data
structure that, once established for one pixel, can be easily updated
to represent the kernel for a neighboring pixel. The data structure
required to perform the equivalent function for mean and variance
operations is much simpler -- all that is necessary is the sum of
pixel values under the kernel (or sum of squared values for variance
calculations), and the kernel size. The sum can be updated by adding
the new values and subtracting the old and the output computed by
dividing by the kernel size.

A more efficient strategy was proposed in \cite{Crow84} for
rectangular kernels. This approach first computes an accumulator image
in which each $i,j$ is replaced by the sum of voxels ``left'' and
``above'' it. This accumulator image may be computed recursively using
local neighborhood values. The mean of rectangles of any size may then
be computed using accumulated values at the rectangle corners. Figures
\ref{fig:accum} and \ref{fig:accumMean} illustrate the 2D case and the
update formulas are provided in Equations \ref{eq:accum} and
\ref{eq:accumMean}.

\begin{eqnarray}
\label{eq:accum}
b_{i,j} & = & \sum_{x \leq i, y \leq j} a_{x,y} \\ \nonumber
	& = & b_{i,j-1} + b_{i-1, j} + a_{i,j} - b_{i-1, j-1}
\end{eqnarray}
where $b_{i,j}$ is the accumulator value at location $i,j$ and $a_{i,j}$ is the input image intensity at $i,j$.

\begin{equation}
\label{eq:accumMean}
m_{i,j} = \frac{1}{wh} (b_{i+w/2,j+h/2} + b_{i-w/2,j-h/2} - b_{i-w/2,j+h/2} + b_{i+w/2,j-h/2})
\end{equation}
where $m_{i,j}$ is the output of the mean filter at location $i,j$.

\begin{figure}[htbp]
\centering
\includegraphics[scale=0.5]{accum2D}
\caption{Recursive computation of the accumulator image - value at $(i,j)$ computed using Equation \ref{eq:accum}.\label{fig:accum}}
\end{figure}

\begin{figure}[htbp]
\centering
\includegraphics[scale=0.5]{accumMean2d}
\caption{Computation of mean using the accumulated image - value at $(i,j)$ computed using Equation \ref{eq:accumMean}.\label{fig:accumMean}}
\end{figure}

This approach to convolution has been used recently in face detection
applications \cite{Viola2004}, where rectangular filters of many
different sizes and shapes were needed.

\section{ITK implementation}
The filters discusses in this article are contained in 3 packages --
{\em consolidatedMorphology}, {\em fastRankMean} and {\em
boxConvolution}. The most recent versions of these packages can be
obtained from {\bf link to Gaetan's darcs repo}. All of the packages
include multithreaded filters and include many examples.

\subsection{Morphology filters -- consolidatedMorphology}
An earlier version of the {\em consolidatedMorphology} package has
been submitted to the InsightJournal. However many improvements and
bug fixes have been included since then.

The morphology filters utilize a new class -- {\em
FlatStructuringElement} to describe, arbitary and decomposable
structuring elements. It provides the following methods to create
structuring elements with different characteristics:
\begin{itemize}
  \item \verb$Ball()$ generate a ball structuring element. It takes the 
	radius of the structuring element as parameter.
  \item \verb$Box()$ produces a box structuring element. It takes the radius 
	of the structuring element as parameter.
  \item \verb$Poly()$ produces an approximation to a circle or sphere using 
	line structuring elements, taking the desired radius and number of 
	lines as parameters. In the 3D case an approximation of a sphere will 
	be attempted. Caution is advised because this is not well tested yet.
  \item \verb$FromImage()$ produces a structuring element from an image. The image is passed
as parameter to the method. A optional parameter can also be passed: the pixel value
to be considered as the foreground value in the image. This value defaults to
\verb$NumericTraits< PixelType >::max()$. This lets the user visualize the structuring element.
\end{itemize}

The algorithm used by a filter is set by the {\em SetAlgorithm}
method, with options of {\em BASIC}, {\em HISTO}, {\em VHGW} and {\em
ANCHOR} for the traditional ITK implementation, the moving histogram
algorithm, the van Herk, Gil, Werman algorithm and the anchor
algorithm respectively. The VHGW and ANCHOR algorithms can only be
used for Box and Poly structuring elements. VHGW is more thoroughly
tested.

The availability of line morphology operations means that new
operations, such as ``union of openings'' can be implemented easily.

\subsection{Rank and mean filters -- fastRankMean}
The {\em fastRankMean} packge uses the sliding window method to
implement rank and mean filters for arbitary kernel shapes. It also
uses line versions of these kernels to provide a ``separable rank''
filter -- an fast approximation to a real rank filter as discussed in
Section \ref{sect:rank}. There are also a number of filters that
accept a mask input and return the rank or mean value of the
intersection of kernel and mask. These are experimental and haven't
been thoroughly tested.

\subsection{Rectangular convolution -- boxConvolution}
The {\em boxConvolution} uses the accumulation method described in
Section \ref{sect:meanVar} to implement rectangular mean and variance
computations. The filters implement different boundary conditions to
{\em itkMeanImageFilter}, with the mean of the kernel inside the image
being computed. The filters in this package are multithreaded and deal
with arbitary dimensions.

\section{Performance}
\subsection{Morpholgy}


\subsection{Rank and Mean filters}

\subsection{Box Convolution}
The relative times for optimized and standard convolution are shown in
Table \ref{tab:perfBoxConv}, and demonstrate significantly improved
performance for kernel radii greater than 2. The times indicate the
theoretical prediction that the complexity of the algorithm is
independent of kernel size.
\begin{table}[htbp]
\centering
\begin{tabular}{ccc}
\hline
Size  &  Direct &  Box \\
1  &     0.014205    &    0.0256313 \\
2  &     0.0355665   &    0.0257389 \\
3  &     0.0507895   &    0.0194342 \\
4  &     0.0745054   &    0.0152253 \\
5  &     0.124225    &    0.0149731 \\
6  &     0.187713    &    0.0149678 \\
7  &     0.27351     &    0.0152494 \\
8  &     0.411708    &    0.0155574 \\
9  &     0.518938    &    0.0156469 \\
10 &     0.682868    &    0.0158462 \\
15  &    2.02712     &    0.0165146 \\
20  &    4.46925     &    0.0176763 \\
25  &    8.43348      & 0.0182798 \\
\hline
\hline
\end{tabular}
\caption{Execution times for direct and box convolution implementations applied to the $256 \times 256$ cthead image on a ??.\label{tab:perfBoxConv}}
\end{table}

\section{Comments and Conclusions}

\section{Acknowledgments}
We thank Dr Pierre Adenot and MIMA2 confocal facilities
(\url{http://mima2.jouy.inra.fr}) for providing the 3D test image.



\appendix



\bibliographystyle{plain}
\bibliography{InsightJournal,local}
\nocite{ITKSoftwareGuide}

\end{document}

